%% Closing chapter: open problems and the future (new for the book)
\chapter{Open Problems and the Future}\label{ch:future}
\chaptermark{Open Problems and the Future}

\begin{chaptersummary}
Each chapter of this book closes with the problems it leaves open. This chapter collects them,
arranges them by the three themes of the book, and states as many as possible as precise
mathematical or statistical questions. It then adds the larger questions that the genomic era is
now forcing on statistical genetics: samples comparable in size to the population, genealogies
inferred at biobank scale, ancient DNA as a time series, the portability of polygenic prediction,
and learned inference whose calibration has to be checked rather than assumed. Short calculations show
the scale of the difficulties: the slow approach of the singleton fraction to Pitman's discount
parameter, the gap between marginal and conditional calibration, and the failure of the
coalescent when the sample approaches the size of the population.
\end{chaptersummary}

\section{What a century has settled}

It is worth beginning with what is no longer in question. Particulate inheritance conserves
variance, and that answers Jenkin. The correlations between relatives follow from Mendel's laws
applied to many loci, and that answers the biometricians. The Fundamental Theorem, read as Fisher
almost certainly meant it and as Ewens has set out, is an identity about the partial change in mean
fitness attributable to changes in gene frequency, and Chapter~\ref{ch:ftns} shows exactly what the
total change adds to it. The Ewens sampling formula is the law of an exchangeable partition whose
predictive rule is affine in the empirical distribution function, and Chapter~\ref{ch:esf} shows
that exchangeability then forces everything else. Under the infinite alleles model the number of
alleles is sufficient for the mutation parameter, and under the infinite sites model the number of
segregating sites is not (Chapters~\ref{ch:suff} and \ref{ch:coal}).

What remains open is, almost without exception, statistical. The mathematics of the classical
models is in excellent shape. What we lack is a reliable account of what samples of the size and
kind we now collect can tell us about the quantities the classical theory defines. Below, each problem is described in words and then, where it is sharp enough, stated
as a displayed question or conjecture precise enough that one could recognise a solution.

\section{Selection and the additive model}

This section collects the problems left by Part I. They concern the empirical content of the
Fundamental Theorem, the prior on genetic architecture, and the estimation of quadratic functionals
of many small effects.

\paragraph{The measurement that has not been made.} The neglected term in the Fundamental Theorem
vanishes exactly when the population moves tangentially to the manifold of Hardy Weinberg and
linkage equilibrium (Chapter~\ref{ch:ftns}). Whether real populations do so on the timescale of
selection is an empirical question, and it is answerable: one needs Lindley's Hardy Weinberg
disequilibrium coefficient and the linkage disequilibrium spectrum, tracked through time in a
population under measurable selection. Long running evolution experiments and the temporal
samples that ancient DNA now provides are the natural places to look. If the association
coordinates turn out to be stationary, the constant genic environment becomes an observation
rather than an assumption. If they move, the size of the remainder can be estimated directly. Of
all the problems in this chapter, this is the one most worth doing first.

Write the total rate of change of mean fitness as $\sigma_A^2(t)+R(t)$, where $\sigma_A^2(t)$ is
the additive genetic variance in the Malthusian parameter and $R(t)$ is the remainder of Chapter~\ref{ch:ftns}, the term that vanishes exactly when the motion is tangential to the manifold of product measures.

\begin{question}\label{future:q:remainder}
Given genotype samples of sizes $n_0,\dots,n_m$ taken at times $t_0<\dots<t_m$ from a population
under selection, construct an estimator of the ratio
\begin{equation*}
\rho=\frac{\int_{t_0}^{t_m}R(t)\dd t}{\int_{t_0}^{t_m}\sigma_A^2(t)\dd t}
\end{equation*}
that is consistent as $\min_i n_i\to\infty$, and determine the sample sizes and spacing needed to
distinguish $\rho=0$ from $|\rho|\ge\rho_0$ with given error probabilities.
\end{question}

The difficulty is one of scale. At a single biallelic locus the departure from Hardy Weinberg
proportions is measured by $F$, and the chi square test has statistic $N\hat F^2$ (Exercise~\ref{intro:exr:F} of
Chapter~\ref{ch:intro}). A one degree of freedom test at level $0.05$ has power $0.8$ when the
noncentrality is $7.85$, so detecting $F=0.1$ needs about $785$ genotypes, detecting $F=0.03$
about $8700$, and detecting $F=0.01$ about $78{,}500$. Departures produced by selection within a
generation are of the order of the selection coefficients, so the measurement requires either
strong selection or samples of biobank size, repeated through time.

\paragraph{Is the genome sparse?} Every estimate in Part I depends on the answer. A horseshoe prior
encodes a genetic architecture with a few substantial effects among many negligible ones; a
Gaussian prior encodes Fisher's infinitesimal model; the omnigenic model of Boyle, Li and Pritchard
suggests that the truth may be dense in a way neither captures. Chapter~\ref{ch:mp} shows, through an exact identity of marginal laws, that the
choice of prior corresponds to the choice of Pitman's discount $\sigma$ (written $\alpha$ in Chapter~\ref{ch:mp}), and that the growth of the
singleton count identifies it. That makes the sparsity level estimable from the frequency
spectrum rather than chosen, but the identification has so far been shown only for the matched
species sampling model.

Recall the correspondence of Chapter~\ref{ch:mp}. If the shrinkage weight of a prior in the three
parameter beta normal class is $\kappa\sim\mathrm{Be}(a,b)$, then $\kappa$ has the law of the
first stick breaking weight of $\mathrm{GEM}(\sigma,\vartheta)$ with $\sigma=1-a$ and
$\vartheta=a+b-1$, and the frequency spectrum intensities agree, $x^{-1}f_\kappa(x)\propto
x^{-1-\sigma}(1-x)^{\vartheta+\sigma-1}$. The horseshoe, $a=b=\tfrac12$, corresponds to
$\sigma=\tfrac12$ and $\vartheta=0$.

\begin{question}\label{future:q:mp}
Let $\Pi_n$ be the exchangeable random partition of $\{1,\dots,n\}$ induced by the global local
shrinkage prior with $\kappa\sim\mathrm{Be}(a,b)$, constructed as in Chapter~\ref{ch:mp}. Is the
exchangeable partition probability function of $\Pi_n$ equal, for every $n$, to Pitman's
\begin{equation*}
p(n_1,\dots,n_k)=\frac{\prod_{i=1}^{k-1}(\vartheta+i\sigma)}{(\vartheta+1)_{n-1}}
\prod_{j=1}^{k}(1-\sigma)_{n_j-1},
\qquad \sigma=1-a,\ \vartheta=a+b-1,
\end{equation*}
where $(x)_m=x(x+1)\cdots(x+m-1)$? If not, which functionals of the partition (beyond the
expected frequency spectrum) do the two laws share?
\end{question}

A positive answer would make the sparsity level a parameter with a sufficient statistic in the
data, rather than a modelling choice. Under the two parameter law the proportion of alleles seen
exactly $j$ times tends to $\sigma(1-\sigma)_{j-1}/j!$, so for the horseshoe half of all alleles
in a large sample are singletons, an eighth are doubletons and a sixteenth tripletons. Example~\ref{future:ex:singleton} shows how slowly this limit is
approached.

\paragraph{Functionals, not coefficients.} The additive genetic variance is a quadratic functional
of the average effects, Robertson's secondary theorem a bilinear one, and under linkage
disequilibrium the estimand is $\alpha^\top V\alpha$ rather than a diagonal form
(Chapter~\ref{ch:hs}). Posterior intervals for such functionals under shrinkage priors undercover
by a few points at moderate sample sizes. The theory for signals of intermediate size suggests that the misses are one sided and that the interval sits low, but the simulations of Chapter~\ref{ch:hs} show the opposite for the plain horseshoe: its misses lie mostly above the truth (80 of 91 in 600 replicates). A frequentist
theory of coverage for quadratic functionals under the horseshoe, analogous to the results of
\citet{future:vanderpas2017} for the coefficients themselves, would make the rate of adaptation a
quantity one could report with a calibrated interval.

\begin{question}\label{future:q:quad}
In the normal means model $Y_i\sim N(\beta_i,1)$, $i=1,\dots,p$, with $\beta$ in the nearly black
class $\ell_0[s_p]=\{\beta:\#\{i:\beta_i\ne0\}\le s_p\}$ and $s_p=o(p)$, let $C_p$ be the
equal tailed $1-\gamma$ horseshoe credible interval for $Q(\beta)=\sum_i\beta_i^2$. Characterise
the set of $\beta\in\ell_0[s_p]$ on which $\liminf_p\Prob_\beta\{Q(\beta)\in C_p\}\ge1-\gamma$,
and decide on which side of $Q(\beta)$ the interval lies, with probability tending to
one, on its complement. The theory of intermediate signals suggests that it lies below, whereas the simulations of Chapter~\ref{ch:hs} found the misses mostly above the truth for the plain horseshoe, in a design that need not lie in the regime of the theory. Both the set to be characterised and the direction on its complement are open, and the second presupposes that the complement has no sequences with limiting coverage strictly between $0$ and $1-\gamma$.
\end{question}

The need for shrinkage, and the danger of it, are both visible in a small calculation. Take
$p=10{,}000$ coefficients of which $100$ equal $3$ and the rest are zero, so that $Q=900$. The plug
in estimate $\sum_iY_i^2$ has expectation $Q+p=10{,}900$, twelve times too large. The unbiased
estimate $\sum_i(Y_i^2-1)$ has standard deviation $(2p+4Q)^{1/2}=154$, so it is correct on average
but noisy. A shrinkage estimator removes most of the noise, and the open question is how much of
the signal it removes with it.

\paragraph{Ascertainment is not a shrinkage problem.} Measuring fitness only on survivors
understates the additive variance severalfold, and no prior repairs it, because the bias is in the
estimand rather than in the estimator. Writing the selection mechanism into the likelihood,
in the manner of the weighted distributions of Rao that Part II applies to the paintbox, is the
obvious route. Chapter~\ref{ch:hs} simulates truncation at a known fraction and rescales, but writing the mechanism into the likelihood has not been carried out for quantitative genetic designs, and it is the
precondition for any empirical test of the Fundamental Theorem in a wild population.

\begin{question}\label{future:q:ascert}
Let $(G,W)$ be the genotype and fitness of an individual, and suppose $(G,W)$ is observed only
with probability $\pi(W)$ for a known or parametrically specified function $\pi$. Under what
conditions on $\pi$ is the additive genetic variance in $W$ identified from the law of the
observed pairs, and what is the efficiency loss relative to complete observation?
\end{question}

\paragraph{Linkage disequilibrium at biobank scale.} The inflation of the plug in additive variance
is $\sigma^2\,\mathrm{tr}(H^\top H)/N$, where $H$ is the smoother matrix of Chapter~\ref{ch:mp}. The trace equals $p$ for the unregularised projection onto $p$ linearly independent columns, whatever the design, and is smaller for a shrinkage smoother.
With real linkage disequilibrium the trace must be computed from an empirical Gram matrix with
strong local correlation, and the effective number of parameters is smaller than $p$ by an amount
that depends on the population's history. Relating that effective dimension to the inferred
genealogy of the sample, which tree sequence methods now make available, would connect Part I to
Part III.

\begin{question}\label{future:q:ld}
For genotype matrices generated by the coalescent with recombination, with scaled recombination
rate $\varrho$ per unit length and $p$ sites per unit length, how does the effective dimension
$p_{\mathrm{eff}}=\mathrm{tr}(H^\top H)$ of Chapter~\ref{ch:mp} behave as $p\to\infty$ and
$N\to\infty$, and can it be estimated from an inferred tree sequence without forming the Gram
matrix?
\end{question}

\section{The sampling theory of neutral alleles}

The problems left by Part II are about extending exact results beyond the families in which they
were proved, and about second order theory for the rare variant end of the frequency spectrum.

\paragraph{Ascertainment and pool size together.} Chapter~\ref{ch:esf} applies two distinct
operations to Kingman's paintbox: a Rao weighting for ascertainment, and the two
parameter family with negative discount for a finite pool of types. Nothing forbids applying both. The resulting three
parameter model would allow ascertainment effort and pool size to be estimated jointly, but its
partition probability function has no closed form, so inference must be simulation based, which is
exactly the setting of Part III.

\begin{question}\label{future:q:three}
Is the three parameter model obtained by Rao weighting the two parameter paintbox identifiable from
the partition of a single sample as $n\to\infty$, and if so at what rate can each parameter be
estimated?
\end{question}

\paragraph{Rigidity beyond tail sum weights.} Chapter~\ref{ch:py} proves that no sampling density
acting on the tail sums of the stick breaking construction preserves the Ewens Pitman family when
$\sigma>0$. Ascertainment schemes that depend on the atom index as well as its position are not
covered. Whether any of them that charges every interval preserves the family is open, and a negative answer would make the
rigidity theorem complete.

\begin{conjecture}\label{future:conj:rigid}
Let $0<\sigma<1$ and $\vartheta>-\sigma$. There is no sampling weight, depending measurably on the
atom index and the atom's position in the subdivision and charging every interval, that carries $\mathrm{PD}(\sigma,\vartheta)$
to $\mathrm{PD}(\sigma',\vartheta')$ for any $(\sigma',\vartheta')\ne(\sigma,\vartheta)$.
\end{conjecture}

\paragraph{Second order theory for the spectrum.} The singleton
fraction estimates the discount parameter with no normalisation, and its bias under a
deterministic envelope is negligible. The same calculation with the random frequencies in place
of the envelope, carried to the order at which the Gamma fluctuations and the zeta function
correction compete, has not been done. Nor has the second order scale matching between the
abundance based estimator and Hill's estimator of the tail index. These are the calculations that
would turn rules of thumb into procedures for the rare variant end of the spectrum, which is where
the information in very large samples lives.

\begin{example}[How fast does the singleton fraction converge?]\label{future:ex:singleton}
Under the two parameter Chinese restaurant process with discount $\sigma$ and concentration
$\vartheta$, the $(n+1)$th gene is a new allele with probability $(\vartheta+\sigma K_n)/(\vartheta+n)$
and joins a given singleton with probability $(1-\sigma)/(\vartheta+n)$. Taking expectations gives
two exact linear recursions,
\begin{align*}
\E K_{n+1}&=\E K_n+\frac{\vartheta+\sigma\E K_n}{\vartheta+n},\\
\E K_{n+1,1}&=\E K_{n,1}\Bigl(1-\frac{1-\sigma}{\vartheta+n}\Bigr)+\frac{\vartheta+\sigma\E K_n}{\vartheta+n},
\end{align*}
where $K_{n,1}$ is the number of singletons. Table~\ref{future:tab:singleton} shows the ratio
$\E K_{n,1}/\E K_n$ for $\sigma=\tfrac12$. With $\vartheta=0$ the ratio is within $0.003$ of
$\sigma$ by $n=100$. With $\vartheta=10$ it is $0.65$ at $n=100$ and still $0.516$ at $n=10^4$;
the excess decays like $c\,n^{-\sigma}$ with $c\approx1.56$, so a sample of about $6\times10^5$ genes is
needed to bring it below $0.002$. The concentration parameter, which is invisible in the first
order theory, controls the bias at every sample size one is likely to meet. A Monte Carlo check
with $4000$ samples of size $100$ at $\vartheta=1$ agrees with the recursions to within
simulation error.
\end{example}

\begin{table}[ht]
\centering
\small
\caption{Ratio $\E K_{n,1}/\E K_n$ of expected singletons to expected alleles under the two
parameter Chinese restaurant process with $\sigma=1/2$. The limit is $\sigma=0.5$ in every
column.}\label{future:tab:singleton}
\begin{tabular}{rccc}
\toprule
$n$ & $\vartheta=0$ & $\vartheta=1$ & $\vartheta=10$ \\
\midrule
$10^2$ & $0.5025$ & $0.5457$ & $0.6506$ \\
$10^3$ & $0.5003$ & $0.5142$ & $0.5493$ \\
$10^4$ & $0.5000$ & $0.5044$ & $0.5156$ \\
$10^5$ & $0.5000$ & $0.5014$ & $0.5049$ \\
$10^6$ & $0.5000$ & $0.5004$ & $0.5016$ \\
\bottomrule
\end{tabular}
\end{table}

\begin{question}\label{future:q:second}
Under $\mathrm{PD}(\sigma,\vartheta)$ find the second order expansion
\begin{equation*}
\begin{gathered}\frac{K_{n,1}}{K_n}=\sigma+n^{-\sigma/2}\,Z_n,\\ \E\frac{K_{n,1}}{K_n}=\sigma+c_1(\sigma,\vartheta)\,n^{-\sigma}+o(n^{-\sigma}),\end{gathered}
\end{equation*}
identifying $c_1$ and the limit law of $Z_n$, and use it to give a bias corrected estimator of
$\sigma$ with an interval whose coverage is correct to second order.
\end{question}

\paragraph{Which functionals identify which parameters?} The de Finetti limit of the two parameter
family is identified by two functionals of the limiting empirical distribution function: the
variance of its heights and the growth rate of its jump count (Chapter~\ref{ch:esf}). The Gibbs
partitions of \citet{future:gnedin2006} are the place to start.

\begin{question}\label{future:q:ident}
For which families of exchangeable partitions is the directing random measure identified by a
finite list of functionals of the limiting empirical distribution function, and is there a
principle that assigns one functional to each parameter?
\end{question}

\paragraph{Priors for concentration parameters.} Chapters~\ref{ch:pig} and \ref{ch:rc} give two
exact routes to Bayesian inference for a concentration parameter: P\'olya inverse gamma
augmentation, which reaches the transcendental factor directly, and rational conjugacy, which works
whenever the Gamma ratio has an integer shift. The two are complementary, and a conjectural meeting point is the half horseshoe, which is a
scale mixture of exponentials (Chapter~\ref{ch:rc} uses it both with residues and, through Proposition~\ref{rc:prop:plugin}, with a latent exponential scale that needs no residues; combining that proposition with the augmentation of Chapter~\ref{ch:pig} has not been carried out here). Two questions remain. The residue expansion of rational conjugacy
is exact but ill conditioned, losing about $n\log_{10}2$ significant digits in a solvable case, or
thirty digits at $n=100$; a stable algorithm for large samples is needed. And the Bernstein and von
Mises argument of Chapter~\ref{ch:pig} shows that for the Ewens richness parameter the Bayes and
classical answers agree to order $1/K$; the interesting cases are therefore the ones where $K$ is
small or where the shape parameter departs from the logseries limit, and these have not been
studied systematically in genetic data.

\begin{question}\label{future:q:stable}
Is there an algorithm that evaluates the normalising constant and posterior moments of a
Pochhammer posterior after $n$ observations with relative error $\epsilon$ in $O(n\log(1/\epsilon))$
operations in fixed precision arithmetic?
\end{question}

\section{Inference without a likelihood}

The problems left by Part III concern what learned summaries can and cannot capture, and how to
tell whether a learned posterior can be trusted.

\paragraph{The ancestral recombination graph.} With recombination the data are a collection of
correlated genealogies and even one dimensional summaries lack closed form laws. The bottleneck
test of Chapter~\ref{ch:suff} applies without change, and its natural output is not another
estimate of the recombination rate but the dimension at which the held out risk stops falling:
an estimate of how many numbers a recombining sample is worth.

For a model $\{P_\theta\}$ for data $X$ define the \emph{bottleneck dimension} $d^*$ as the
smallest $d$ for which some $T:\mathcal X\to\R^d$ from a fixed class of encoders, as in Chapter~\ref{ch:suff}, invariant under the symmetries of the model, is sufficient for $\theta$. The class matters: if $T$ may be any measurable map then $d^*\le1$ for every model on a countable sample space, because the orbit of $X$ under the symmetries is sufficient and can be coded injectively by a single real number.

\begin{question}\label{future:q:arg}
For a sample of $n$ sequences of fixed length under the coalescent with recombination at scaled
rate $\varrho>0$ and infinite sites mutation at scaled rate $\vartheta$, is $d^*$ finite? If it is,
how does it grow with $n$ and with the sequence length?
\end{question}

\paragraph{Selection in the genealogy.} Tajima's $D$ and its relatives are hand built two
dimensional summaries designed to detect departures from the neutral frequency spectrum. A learned
two dimensional bottleneck fitted under models with directional or balancing selection is a direct
test of how much the classical constructions leave on the table.

\paragraph{Multiple merger coalescents.} For species with high fecundity and sweepstakes
reproduction the Kingman coalescent is replaced by $\Lambda$ coalescents, for which the sufficiency
question is genuinely open while simulation remains easy. Here the bottleneck dimension is itself
the scientific quantity, since it counts the independent features of the spectrum that distinguish
a multiple merger history from a neutral one.

\begin{question}\label{future:q:beta}
For the Beta$(2-\alpha,\alpha)$ coalescent with $1<\alpha<2$ and infinite sites mutation, is the
site frequency spectrum sufficient for $(\alpha,\vartheta)$? If not, what is $d^*$?
\end{question}

\paragraph{Beyond the permutation invariant ceiling.} Chapter~\ref{ch:coal} proves that any encoder
built from per site permutation invariant features factors through the site frequency spectrum, and
that the number of distinct haplotypes is a minimal witness that escapes the ceiling. A small
case shows what the spectrum loses. With $n=5$ sequences and two segregating sites, each carried
by two sequences, the spectrum is the same whether both mutations sit on the same pair of
sequences (haplotypes $11,11,00,00,00$, so two distinct haplotypes) or on disjoint pairs
($10,10,01,01,00$, three distinct haplotypes). Both are consistent with a single tree, but they
require different trees, and they have different probabilities under the coalescent.

\begin{question}\label{future:q:minsuff}
Under the infinite sites model without recombination, with the sample represented as an $n\times S$
binary matrix modulo permutations of rows and of columns, what is the minimal sufficient statistic
for $\vartheta$? Which architectures, built from equivariant layers, can represent it?
\end{question}

\paragraph{Calibration, not just accuracy.} Learned posteriors can match exact posteriors to within
a few percent and still fail simulation based calibration, and Chapter~\ref{ch:coal} shows that
standard rank recalibration removes the marginal error but cannot show that the conditional error is gone, and
a mild conditional residual remains. A theory of conditional calibration for amortised posteriors, with
diagnostics that a practitioner can run, is perhaps the most important methodological problem in
likelihood free inference today. The next example shows why marginal checks cannot be enough.

\begin{example}[A marginally calibrated posterior that is conditionally wrong]\label{future:ex:calib}
Let $\theta\sim N(0,1)$ and $X\mid\theta\sim N(\theta,1)$, so the exact posterior is
$N(x/2,1/2)$. Suppose a learned method ignores the data and reports the prior $N(0,1)$ as its
posterior. Simulation based calibration \citep{future:cook2006} draws $\theta$ from the prior,
data from the model, and computes the rank of $\theta$ in the reported posterior; here the rank
statistic is $\Phi(\theta)$, which is exactly uniform. The method passes the check perfectly, and
its $90$ percent interval $(-1.645,1.645)$ covers $\theta$ with marginal probability exactly
$0.90$. Conditionally it fails badly. Given $X=x$ the coverage is
$\Prob(|\theta|<1.645\mid x)$ with $\theta\sim N(x/2,1/2)$, which is $0.98$ at $x=0$, $0.82$ at
$x=2$ and $0.58$ at $x=3$; over the $16$ percent of data sets with $|X|>2$ a simulation gives
coverage $0.67$. Overcoverage near the centre of the data distribution exactly pays for
undercoverage in the tails.
\end{example}

\begin{question}\label{future:q:calib}
Let $q(\theta\mid x)$ be an amortised posterior and $U=Q(\theta\mid X)$ the probability integral
transform of the true parameter under it. Simulation based calibration tests whether $U$ is
uniform. Find a test, computable from simulated pairs $(\theta_i,X_i)$ alone, that is consistent
against every alternative in which $U$ is not uniform conditionally on $X$ on a set of positive
probability, and give its rate against alternatives at distance $\delta$ in a suitable metric.
\end{question}

\paragraph{Misspecification.} Every method in Part III conditions on the simulator being correct.
When it is not, a learned posterior is confidently wrong, arguably more so than approximate Bayesian
computation, whose tolerance inadvertently widens the answer. Posterior predictive checks against
the observed spectrum are a start. Generalised Bayesian updating and the cut models of modular
inference offer more principled options, and their application to coalescent models has hardly
begun.

\section{The larger questions}

The problems above are local to the chapters of this book. Several larger ones are reshaping the
whole field, and they are less sharply posed; for most of them the first task is to find the right
question.

\paragraph{Samples as large as populations.} The coalescent is an approximation for a small sample
from a large population. When the sample is comparable to the effective population size, as it now
is for some human biobanks, multiple and simultaneous mergers appear in the recent past, the
frequency spectrum at the rare end departs from the Kingman prediction, and estimates of recent
demography and of selection on rare variants are distorted \citep{future:bhaskar2014}. The exact
Wright Fisher genealogy is tractable numerically, and replacing the coalescent by it, or by a
hybrid, is under way.

\begin{example}[The birthday problem in one generation]\label{future:ex:birthday}
Trace $n$ lineages back one generation in a Wright Fisher population of $2N$ genes. Each lineage
picks its parent uniformly and independently, so this is the birthday problem with $2N$ days. The
expected number of pairs that pick the same parent is $\binom n2/(2N)$, and the probability that
no two do is $\prod_{i=1}^{n-1}(1-i/2N)$. With $2N=20{,}000$ and $n=100$ the expected number of
merging pairs is $0.25$ and the probability of no merger is $0.78$, so the Kingman picture of at
most one rare merger per generation is reasonable. With $n=1000$ the expected number is $25$, the
probability of no merger is $9\times10^{-12}$, and the sample loses on average $24.6$ lineages in
the first generation, with simultaneous mergers and occasional triple mergers. With $n=5000$ about
$576$ lineages are lost at once. The Kingman coalescent, which allows exactly one pairwise merger
at a time, describes none of this, and the errors fall precisely on the recent branches that carry
the rare variants.
\end{example}

\begin{question}\label{future:q:finite}
Let a sample of $n$ genes be drawn without replacement from a Wright Fisher population of $2N$
genes at stationarity under the infinite alleles model, with $n/(2N)\to c\in(0,1]$ and
$4N\mu\to\vartheta$. What is the limiting law of the allelic partition of the sample, and how does
it differ from the Ewens sampling formula in the counts of rare alleles?
\end{question}

\paragraph{Genealogies as data.} Methods that infer the ancestral recombination graph of hundreds of
thousands of genomes have turned the genealogy from a latent variable into something close to an
observation. This changes the statistical questions. Inference about mutation, recombination,
demography and selection can be conditioned on an estimated genealogy, but the uncertainty in that
estimate has to be propagated, and the estimated genealogies are themselves the output of
approximate algorithms whose biases are not well characterised. A Bayesian treatment that is both
coherent and feasible at this scale does not yet exist.

\paragraph{Time series of genomes.} Ancient DNA has given population genetics what it has rarely had:
direct observation of allele frequencies through time. Selection can then be measured as a change in
frequency rather than inferred from a pattern of variation at one time. But ancient samples are
small, degraded, unevenly spaced, drawn from shifting populations, and subject to strong
ascertainment. They are, in the language of Part II, a sample drawn with a nonuniform density from
a moving paintbox, and the theory of Chapters~\ref{ch:esf} and \ref{ch:py} is a beginning, not an
answer.

\paragraph{Polygenic adaptation and confounding.} Signals of polygenic adaptation for human height,
detected by combining many small allele frequency differences, were substantially reduced when
recomputed with association effect sizes estimated in a more homogeneous sample
\citep{future:berg2019,future:sohail2019}. The lesson is that subtle population stratification in
the estimated effects aggregates, exactly as the inflation of Chapter~\ref{ch:mp} aggregates, into
a spurious signal. The methodological problem is to estimate a many parameter functional when the
errors in the parameters are correlated with the very structure the functional is meant to detect.

\paragraph{Portability of prediction.} Polygenic scores built from samples of predominantly
European ancestry predict much less accurately in other populations \citep{future:martin2019}. Part
of the loss is linkage disequilibrium, part is differences in allele frequencies, and part may be
genuine differences in effect or in environment. Separating these is a problem in transfer learning
with a known genealogical structure, and the hierarchical models of Part I, with population
specific local scales sharing a global scale, are a natural starting point.

\paragraph{Learned models of the genome.} Large neural networks trained on genomic sequence are now
used to predict the functional effects of variants. For population genetics they are potentially a
source of prior information: an informative prior on which loci are likely to matter, which is
precisely what the local scales of a global local prior need. Using them responsibly requires
treating their outputs as priors to be updated, not as answers, and checking the calibration of the
resulting posteriors. This is the same discipline that Part III asks of learned posteriors, now
applied to learned priors.

\section{A closing note}

Fisher regarded the Fundamental Theorem as one of his deepest results, and Ewens spent much of a
long career showing what it does and does not say. The sampling formula that Ewens derived in 1972
became, through Kingman, Ferguson, Antoniak and Pitman, one of the central objects of modern
probability and machine learning, far from the electrophoresis gels that motivated it. Both stories
show the same pattern: a statement about genes becomes a statement about statistics, and the
statistics then returns, years later, to change how the genetics is done.

The thread that runs through this book is that the questions of population genetics are, at bottom,
questions of inference: what a sample can tell us about a population, what a prior commits us to,
and what is lost when we replace a likelihood by a summary or a simulation. Fisher, Haldane, Wright,
Kimura, Ewens and Kingman answered the first of these questions for the models of their day with
remarkable completeness. The data we have now are larger and stranger than any of them imagined,
and the questions are open again.

\section{Exercises}

\begin{enumerate}[label=\arabic*.,ref=\arabic*,leftmargin=*]

\item Show that a one degree of freedom chi square test at level $0.05$ has power $0.8$ when the
noncentrality parameter is $7.85$, and deduce the sample sizes quoted after
Question~\ref{future:q:remainder}. \emph{Hint.} Under the alternative $N\hat F^2$ is approximately
noncentral chi square with noncentrality $NF^2$.

\item Verify that if $\kappa\sim\mathrm{Be}(a,b)$ then $x^{-1}f_\kappa(x)$ is proportional to the
intensity $x^{-1-\sigma}(1-x)^{\vartheta+\sigma-1}$ of $\mathrm{PD}(\sigma,\vartheta)$ with
$\sigma=1-a$ and $\vartheta=a+b-1$, and that the first stick breaking weight of
$\mathrm{GEM}(\sigma,\vartheta)$ is $\mathrm{Be}(1-\sigma,\vartheta+\sigma)=\mathrm{Be}(a,b)$.
\emph{Hint.} Compare exponents.

\item Show that the limiting proportion of alleles seen $j$ times under $\mathrm{PD}(\sigma,\vartheta)$
is $\sigma(1-\sigma)_{j-1}/j!$, and check that these proportions sum to one. Evaluate them for
$\sigma=\tfrac12$ and $j\le5$. \emph{Hint.} Use the generating function $1-(1-z)^\sigma$.

\item In the normal means example after Question~\ref{future:q:quad}, show that
$\E\sum_iY_i^2=Q+p$ and $\Var\sum_i(Y_i^2-1)=2p+4Q$. \emph{Hint.} $Y_i^2$ is noncentral chi
square with one degree of freedom.

\item Derive the two recursions of Example~\ref{future:ex:singleton} and reproduce
Table~\ref{future:tab:singleton}. Then estimate the constant $c$ of that example numerically for $\vartheta=1$ and $\vartheta=10$ from the ratio of
expectations. \emph{Hint.} $(\E K_{n,1}/\E K_n-\sigma)\,n^{\sigma}$ settles near $0.44$ and $1.56$.

\item In Example~\ref{future:ex:calib} compute the conditional coverage of the exact $90$ percent
posterior interval $x/2\pm1.645/\sqrt2$, and show that it equals $0.90$ for every $x$. Show also
that a reported posterior of the form $N(\lambda x,v)$ passes simulation based calibration if
and only if $v=(1-\lambda)^2+\lambda^2$, so that a whole curve of wrong posteriors passes, with
the prior ($\lambda=0$) at one end and the exact posterior ($\lambda=\tfrac12$) at its minimum.
\emph{Hint.} The rank statistic is $\Phi\{(\theta-\lambda X)/\sqrt v\}$, and
$\theta-\lambda X$ is Gaussian with mean zero and variance $(1-\lambda)^2+\lambda^2$.

\item Using the birthday calculation of Example~\ref{future:ex:birthday}, find the largest $n$ for
which the probability of at least one merger in the first generation back is below $0.5$ when
$2N=20{,}000$, and compare with $\sqrt{2N\cdot2\log2}$. \emph{Hint.} The answer is $166$.

\item In Example~\ref{future:ex:birthday} show that the expected number of distinct parents of the
$n$ lineages is $2N\{1-(1-1/2N)^n\}$, and deduce the losses of $24.6$ lineages at $n=1000$ and
about $576$ at $n=5000$. Show further that the expected loss is
$\binom n2/(2N)-\binom n3/(2N)^2+\cdots$, so that it falls short of the expected number of merging
pairs by about $\binom n3/(2N)^2$, and explain why. \emph{Hint.} Write the number of distinct
parents as a sum over the $2N$ possible parents of the indicator that at least one lineage picks
that parent, and expand $(1-1/2N)^n$ by the binomial theorem; a triple merger counts three pairs
but loses only two lineages.

\item In Example~\ref{future:ex:calib} verify the conditional coverages $0.98$, $0.82$ and $0.58$
of the reported interval $(-1.645,1.645)$ at $x=0,2,3$, show that $\Prob(|X|>2)=0.157$, and
compute the coverage over the data sets with $|X|>2$ by numerical integration, confirming the
simulated value $0.67$. \emph{Hint.} Marginally $X\sim N(0,2)$, and given $X=x$ the coverage is
$\Phi\{\sqrt2(1.645-x/2)\}-\Phi\{\sqrt2(-1.645-x/2)\}$.

\item Prove the remark before Question~\ref{future:q:arg}: if the sample space is countable and
every $P_\theta$ is invariant under a group $G$ acting on it, then the orbit $GX$ is sufficient
for $\theta$, and so is a single real number. \emph{Hint.} Points of one orbit have equal
probability under every $P_\theta$, so an orbit of positive probability is finite and $X$ is
uniform on it given $GX$; the orbits form a countable set, which can be numbered.


\end{enumerate}

\begin{thebibliography}{9}

\bibitem[Berg et~al.(2019)]{future:berg2019}
Berg, J.~J., Harpak, A., Sinnott Armstrong, N., Joergensen, A.~M., Mostafavi, H., Field, Y.,
Boyle, E.~A., Zhang, X., Racimo, F., Pritchard, J.~K. and Coop, G. (2019). Reduced signal for
polygenic adaptation of height in UK Biobank. \emph{eLife} 8, e39725.

\bibitem[Bhaskar et~al.(2014)]{future:bhaskar2014}
Bhaskar, A., Clark, A.~G. and Song, Y.~S. (2014). Distortion of genealogical properties when the
sample is very large. \emph{Proceedings of the National Academy of Sciences} 111, 2385 to 2390.

\bibitem[Cook et~al.(2006)]{future:cook2006}
Cook, S.~R., Gelman, A. and Rubin, D.~B. (2006). Validation of software for Bayesian models using
posterior quantiles. \emph{Journal of Computational and Graphical Statistics} 15, 675 to 692.

\bibitem[Gnedin and Pitman(2006)]{future:gnedin2006}
Gnedin, A. and Pitman, J. (2006). Exchangeable Gibbs partitions and Stirling triangles.
\emph{Journal of Mathematical Sciences} 138, 5674 to 5685.

\bibitem[Martin et~al.(2019)]{future:martin2019}
Martin, A.~R., Kanai, M., Kamatani, Y., Okada, Y., Neale, B.~M. and Daly, M.~J. (2019). Clinical
use of current polygenic risk scores may exacerbate health disparities. \emph{Nature Genetics} 51,
584 to 591.

\bibitem[Sohail et~al.(2019)]{future:sohail2019}
Sohail, M., Maier, R.~M., Ganna, A., Bloemendal, A., Martin, A.~R., Turchin, M.~C., Chiang, C.~W.~K., Hirschhorn, J., Daly, M.~J., Patterson, N., Neale, B.~M., Mathieson, I., Reich, D. and Sunyaev, S.~R. (2019). Polygenic adaptation on height is overestimated due to uncorrected
stratification in genome wide association studies. \emph{eLife} 8, e39702.

\bibitem[van der Pas et~al.(2017)]{future:vanderpas2017}
van der Pas, S., Szab\'o, B. and van der Vaart, A. (2017). Uncertainty quantification for the
horseshoe (with discussion). \emph{Bayesian Analysis} 12, 1221 to 1274.

\end{thebibliography}
